\ifdefined\gkver\else\def\gkver{student}\fi
\documentclass[\gkver]{gaokaozhenti}
\begin{document}

\chapter{2007年上海}

\begin{enumerate}
\item
%% number: 1
%% paperName: 2007年上海高考第 1 题 4 分
%% typeId: 3
%% type: 填空题
%% chapter: 磁场
%% point: 磁感应强度
%% method: 无
%% score: 4
%% degree: 850
%% duplicateId: 0
%% body:
\textbf{A．}磁场对放入其中的长为 $l$、电流强度为 $I$、方向与磁场垂直的通电直导线有力的作用，可以用磁感应强度 $B$ 描述磁场的力的性质，磁感应强度的大小 $B=$ \tkanswer[1.8cm]{$\frac{F}{Il}$}，用类似方法描述物质基本性质的物理量还有 \tkanswer[2.2cm]{电场强度}。

\textbf{B．}在磁感应强度为 $B$ 的匀强磁场中，垂直于磁场方向放入一段通电直导线。若任意时刻该导线中有 $n$ 个以速度 $v$ 做定向移动的电荷，每个电荷所带电量为 $q$。则每个电荷所受的洛伦兹力 $f=$ \tkanswer[1.6cm]{$qvB$}，该段直导线所受的安培力 $F=$ \tkanswer[1.8cm]{$nqvB$}。

%% answer:
\jdanswer{
\textbf{A．}$\frac{F}{Il}$，电场强度。

\textbf{B．}$qvB$，$nqvB$。
}

%% memo:
\memoanswer{
\textbf{A．}当通电直导线垂直放在磁场中时，受到的安培力 $F=BIl$，故磁感应强度 $B=\frac{F}{Il}$；用比值来定义的物理量还有电场强度 $E=\frac{F}{q}$。

\textbf{B．}当带电粒子的速度方向与磁场方向垂直时，每个电荷受到的洛伦兹力 $f=qvB$；导线受到的安培力可以看成是 $n$ 个粒子受到的洛伦兹力的合力，$F=nqvB$。
}

\item
%% number: 2
%% paperName: 2007年上海高考第 2 题 4 分
%% typeId: 3
%% type: 填空题
%% chapter: 机械振动
%% point: 受迫振动共振
%% method: 无
%% score: 4
%% degree: 650
%% duplicateId: 0
%% body:
\textbf{A．}沿 $x$ 轴正方向传播的简谐横波在 $t=0$ 时的波形如图所示。P、Q 两个质点的平衡位置分别位于 $x=3.5\Um$ 和 $x=6.5\Um$ 处。在 $t_{1}=0.5\Us$ 时，质点 P 恰好第二次处于波峰位置；则 $t_{2}=$ \tkanswer[1.4cm]{0.6} $\Us$ 时，质点 Q 第二次处在平衡位置且向上运动；当 $t_{3}=0.9\Us$ 时，质点 P 的位移为 \tkanswer[1.4cm]{2} $\Ucm$。

\onepicture[w=7.5cm, align=right]{\includesvg{figs/02}}

\textbf{B．}在接近收费口的道路上安装了若干条突起于路面且与行驶方向垂直的减速带。减速带间距为 $10\Um$，当车辆经过减速带时会产生振动。若某汽车的固有频率为 $1.25\UHz$，则当该车以 \tkanswer[1.6cm]{12.5} $\Ums$ 的速度行驶在此减速区时颠簸得最厉害，我们把这种现象称为 \tkanswer[1.6cm]{共振}。

%% answer:
\jdanswer{
\textbf{A．}0.6，2。

\textbf{B．}12.5，共振。
}

%% memo:
\memoanswer{
\textbf{A．}由波向 $x$ 正方向传播，可以断定 $t=0$ 时质点 P 振动方向向上，Q 点由平衡位置向下振动；$t_{1}$ 后第二次处于波峰位置，则 $\frac{5T}{4}=t_{1}$，得 $T=0.4\Us$，Q 第二次处在平衡位置且向上运动的时间是 $t_{2}=\frac{6T}{4}=0.6\Us$；经过 $0.9\Us$，即 $\frac{9T}{4}$，质点 P 到达正向位移最大处，P 点的位移是 $2\Ucm$。

\textbf{B．}车子在减速带的作用下做周期性的振动，$T=\frac{s}{v}$，当振动周期 $T$ 等于固有周期时，即 $T=\frac{1}{f}=\frac{1}{1.25}\Us$ 时发生共振，$v=\frac{s}{T}=\frac{10}{1/1.25}\Ums=12.5\Ums$。
}

\item
%% number: 3
%% paperName: 2007年上海高考第 3 题 4 分
%% typeId: 3
%% type: 填空题
%% chapter: 电路
%% point: 分压与分流
%% method: 无
%% score: 4
%% degree: 650
%% duplicateId: 0
%% body:
\textbf{A．}如图所示，AB 两端接直流稳压电源，$U_{AB}=100\UV$，$R_{0}=40\UO$，滑动变阻器总电阻 $R=20\UO$，当滑动片处于滑动变阻器中点时，C、D 两端电压 $U_{CD}$ 为 \tkanswer[1.4cm]{80} $\UV$，通过电阻 $R_{0}$ 的电流为 \tkanswer[1.4cm]{2} $\UA$。

\onepicture[w=4.2cm, align=right]{\includesvg{figs/03a}}

\textbf{B．}如图所示，自耦变压器输入端 A、B 接交流稳压电源，其电压有效值 $U_{AB}=100\UV$，$R_{0}=40\UO$，当滑动片处于线圈中点位置时，C、D 两端电压的有效值 $U_{CD}$ 为 \tkanswer[1.5cm]{200} $\UV$，通过电阻 $R_{0}$ 的电流的有效值为 \tkanswer[1.4cm]{5} $\UA$。

\onepicture[w=3.8cm, align=right]{\includesvg{figs/03b}}

%% answer:
\jdanswer{
\textbf{A．}80，2。

\textbf{B．}200，5。
}

%% memo:
\memoanswer{
\textbf{A．}$R$ 的上半部分与 $R_{0}$ 串联，总电阻是 $50\UO$，这部分电阻两端的总电压等于电源电压 $100\UV$，由电压分配规律，CD 两端电压是 $80\UV$，通过的电流 $I=\frac{U_{CD}}{R_{0}}=2\UA$。

\textbf{B．}由变压器的性质 $\frac{U_{1}}{U_{2}}=\frac{n_{1}}{n_{2}}$ 得，CD 两端的电压是 $200\UV$，通过电阻的电流 $I=\frac{U_{2}}{R_{0}}=5\UA$。
}

\item
%% number: 4
%% paperName: 2007年上海高考第 4 题 4 分
%% typeId: 3
%% type: 填空题
%% chapter: 原子和原子核
%% point: 天然放射性
%% method: 无
%% score: 4
%% degree: 650
%% duplicateId: 0
%% body:
一置于铅盒中的放射源发射出的 $\alpha$、$\beta$ 和 $\gamma$ 射线，由铅盒的小孔射出，在小孔外放一铝箔，铝箔后的空间有一匀强电场。进入电场后，射线变为 a、b 两束，射线 a 沿原来方向行进，射线 b 发生了偏转，如图所示。则图中的射线 a 为 \tkanswer[1.4cm]{$\gamma$} 射线，射线 b 为 \tkanswer[1.4cm]{$\beta$} 射线。

\onepicture[w=6.5cm, align=right]{\includesvg{figs/04}}

%% answer:
\jdanswer{$\gamma$，$\beta$}

%% memo:
\memoanswer{
$\alpha$ 射线穿透能力最弱，会被铝箔阻挡，无法进入电场，$\beta$ 和 $\gamma$ 射线会进入电场；而 $\gamma$ 射线不带电，不会在电场中发生偏转，$\beta$ 射线带负电，会在电场中发生偏转。因此射线 a 是 $\gamma$ 射线，射线 b 是 $\beta$ 射线。
}

\item
%% number: 5
%% paperName: 2007年上海高考第 5 题 4 分
%% typeId: 3
%% type: 填空题
%% chapter: 功和能
%% point: 机械能守恒定律
%% method: 无
%% score: 4
%% degree: 450
%% duplicateId: 0
%% body:
在竖直平面内，一根光滑金属杆弯成如图所示形状，相应的曲线方程为 $y=2.5\cos\left(kx+\frac{2}{3}\pi\right)$（单位：$\Um$），式中 $k=1\Um^{-1}$。将一光滑小环套在该金属杆上，并从 $x=0$ 处以 $v_{0}=5\Ums$ 的初速度沿杆向下运动，取重力加速度 $g=10\Umsq$。则当小环运动到 $x=\frac{\pi}{3}\Um$ 时的速度大小 $v=$ \tkanswer[1.6cm]{$5\sqrt{2}$} $\Ums$；该小环在 $x$ 轴方向最远能运动到 $x=$ \tkanswer[1.6cm]{$\frac{5}{6}\pi$} $\Um$ 处。

\onepicture[w=6.5cm, align=right]{\includesvg{figs/05}}

%% answer:
\jdanswer{$5\sqrt{2}$，$\frac{5}{6}\pi$}

%% memo:
\memoanswer{
由曲线方程知，环在 $x=0$ 处的 $y$ 坐标是 $-\frac{2.5}{2}\Um$；在 $x=\frac{\pi}{3}$ 时，$y=2.5\cos\left(kx+\frac{2}{3}\pi\right)=-2.5\Um$。由机械能守恒得 $\frac{1}{2}mv_{2}^{2}=\frac{1}{2}mv_{1}^{2}+mg\left(2.5-\frac{2.5}{2}\right)$，则环到达 $x=\frac{\pi}{3}$ 处的速度是 $v_{2}=5\sqrt{2}\Ums$；当环到达最高点时速度为零，由机械能守恒 $0=\frac{1}{2}mv_{1}^{2}+mg\left(-h-\frac{2.5}{2}\right)$，得 $h=0$，即 $y=0$，故该小环在 $x$ 轴方向最远能运动到 $x=\frac{5\pi}{6}\Um$。
}

\item
%% number: 6
%% paperName: 2007年上海高考第 6 题 5 分
%% typeId: 1
%% type: 单选题
%% chapter: 原子和原子核
%% point: 天然放射性
%% method: 无
%% score: 5
%% degree: 850
%% duplicateId: 0
%% body:
$\ce{^{238}_{92}U}$ 衰变为 $\ce{^{222}_{86}Ra}$ 要经过 $m$ 次 $\alpha$ 衰变和 $n$ 次 $\beta$ 衰变，则 $m$，$n$ 分别为\xzanswer{B}

\fourchoices[answer=B]
{2，4}
{4，2}
{4，6}
{16，6}

%% answer: B
%% memo:
\memoanswer{
由于衰变后的新核质量数减少了 16，说明发生 $\alpha$ 衰变的次数是 $m=\frac{16}{4}=4$ 次，发生 $\beta$ 衰变的次数是 $n=86-92+8=2$，故 B 正确。
}

\item
%% number: 7
%% paperName: 2007年上海高考第 7 题 5 分
%% typeId: 1
%% type: 单选题
%% chapter: 磁场
%% point: 磁场的叠加
%% method: 无
%% score: 5
%% degree: 650
%% duplicateId: 0
%% body:
取两根完全相同的长导线，用其中一根绕成如图（a）所示的螺线管，当在该螺线管中通以电流强度为 $I$ 的电流时，测得螺线管内中部的磁感应强度大小为 $B$；若将另一根长导线对折后绕成如图（b）所示的螺线管，并通以电流强度也为 $I$ 的电流，则在该螺线管内中部的磁感应强度大小为\xzanswer{A}

\twopicture[labelA=2007上海07a, labelB=2007上海07b, wA=5.5cm, wB=5.5cm]
{figs/07a.png}
{figs/07b.png}

\fourchoices[answer=A]
{0}
{$0.5B$}
{$B$}
{$2B$}

%% answer: A
%% memo:
\memoanswer{
由于螺线管是来回绕的，所以来回绕行的电流产生的磁场相互抵消，螺线管中的磁感应强度是零，故 A 正确。
}

\item
%% number: 8
%% paperName: 2007年上海高考第 8 题 5 分
%% typeId: 2
%% type: 多选题
%% chapter: 光学
%% point: 光的衍射
%% method: 无
%% score: 5
%% degree: 650
%% duplicateId: 0
%% body:
光通过各种不同的障碍物后会产生各种不同的衍射条纹，衍射条纹的图样与障碍物的形状相对应，这一现象说明\xzanswer{BC}

\fourchoices[answer=BC]
{光是电磁波}
{光具有波动性}
{光可以携带信息}
{光具有波粒二象性}

%% answer: BC
%% memo:
\memoanswer{
光能发生衍射现象，说明光有波动性，故 B 正确；衍射图样与障碍物的形状对应，说明了衍射图样中包含了障碍物等物体的信息，故 C 正确；光是电磁波，光也具有波粒二象性，但在这个现象中没有得到反映，故 AD 错误。
}

\item
%% number: 9
%% paperName: 2007年上海高考第 9 题 5 分
%% typeId: 1
%% type: 单选题
%% chapter: 机械波
%% point: 波长频率波速
%% method: 无
%% score: 5
%% degree: 650
%% duplicateId: 0
%% body:
如图所示，位于介质 Ⅰ 和 Ⅱ 分界面上的波源 S，产生两列分别沿 $x$ 轴负方向与正方向传播的机械波。若在两种介质中波的频率及传播速度分别为 $f_{1}$、$f_{2}$ 和 $v_{1}$、$v_{2}$，则\xzanswer{C}

\onepicture[w=10.0cm]{\includesvg{figs/09}}

\fourchoices[answer=C]
{$f_{1}=2f_{2}$，$v_{1}=v_{2}$}
{$f_{1}=f_{2}$，$v_{1}=0.5v_{2}$}
{$f_{1}=f_{2}$，$v_{1}=2v_{2}$}
{$f_{1}=0.5f_{2}$，$v_{1}=v_{2}$}

%% answer: C
%% memo:
\memoanswer{
由于是同一波源发出的两列波，两列波的频率相同，即 $f_{1}=f_{2}$；由波速公式 $v=\lambda f$，得到两波的波速之比 $\frac{v_{1}}{v_{2}}=\frac{\lambda_{1}}{\lambda_{2}}=\frac{2}{1}$，故 C 正确。
}

\item
%% number: 10
%% paperName: 2007年上海高考第 10 题 5 分
%% typeId: 2
%% type: 多选题
%% chapter: 力物体的平衡
%% point: 三力平衡
%% method: 无
%% score: 5
%% degree: 650
%% duplicateId: 0
%% body:
如图所示，用两根细线把 A、B 两小球悬挂在天花板上的同一点 O，并用第三根细线连接 A、B 两小球，然后用某个力 $F$ 作用在小球 A 上，使三根细线均处于直线状态，且 OB 细线恰好沿竖直方向，两小球均处于静止状态。则该力可能为图中的\xzanswer{BC}

\onepicture[w=4.8cm]{figs/10.png}

\fourchoices[answer=BC]
{$F_{1}$}
{$F_{2}$}
{$F_{3}$}
{$F_{4}$}

%% answer: BC
%% memo:
\memoanswer{
A 小球受三个力的作用：重力 $G$、绳子 OA 向上的拉力 $T$ 和拉力 $F$，绳子 AB 中没有拉力，只有 $G$、$T$、$F$ 三力平衡。由平衡条件知，水平方向和竖直方向的分力都平衡，$F_{4}$ 只有水平向左的分量，不能平衡 $T$ 的水平向左的分量，故 D 错误；$F_{1}$ 竖直分量总比 $T$ 的竖直分量小，所以 $F_{1}$ 水平向右的分量总比 $T$ 的水平分量小，不能平衡，故 A 错误，BC 正确。
}

\item
%% number: 11
%% paperName: 2007年上海高考第 11 题 5 分
%% typeId: 2
%% type: 多选题
%% chapter: 气体性质
%% point: 液柱移动定性分析
%% method: 无
%% score: 5
%% degree: 650
%% duplicateId: 0
%% body:
如图所示，一定质量的空气被水银封闭在静置于竖直平面的 U 型玻璃管内，右管上端开口且足够长，右管内水银面比左管内水银面高 $h$。能使 $h$ 变大的原因是\xzanswer{ACD}

\onepicture[w=3.4cm]{\includesvg{figs/11}}

\fourchoices[answer=ACD]
{环境温度升高}
{大气压强升高}
{沿管壁向右管内加入水银}
{U 型玻璃管自由下落}

%% answer: ACD
%% memo:
\memoanswer{
当环境温度升高时，左端封闭的气体体积要增大，$h$ 增大，故 A 正确；大气压强升高，封闭气体压强增大，体积缩小，$h$ 减小，故 B 错误；沿 U 形管加入水银，封闭气体的压强增大，平衡时内外水银高度差增大，故 C 正确；U 形管自由下落时，水银重力产生加速度而不产生压力，由于原来封闭气体的压强大于大气压强，体积膨胀，$h$ 增大，故 D 正确。
}

\item
%% number: 12
%% paperName: 2007年上海高考第 12 题 5 分
%% typeId: 2
%% type: 多选题
%% chapter: 功和能
%% point: 动能定理
%% method: 无
%% score: 5
%% degree: 650
%% duplicateId: 0
%% body:
物体沿直线运动的 $v-t$ 关系如图所示，已知在第 1 秒内合外力对物体做功为 $W$，则\xzanswer{CD}

\onepicture[w=7.5cm]{\includesvg{figs/12}}

\fourchoices[answer=CD]
{从第 1 秒末到第 3 秒末合外力做功为 $4W$}
{从第 3 秒末到第 5 秒末合外力做功为 $-2W$}
{从第 5 秒末到第 7 秒末合外力做功为 $W$}
{从第 3 秒末到第 4 秒末合外力做功为 $-0.75W$}

%% answer: CD
%% memo:
\memoanswer{
在第 $1\Us$ 内，物体做匀加速运动，由动能定理得合力做的功为 $W=\frac{1}{2}mv^{2}-0$；从第 $1\Us$ 末到第 $3\Us$ 末，合力做的功为 $W_{2}=\frac{1}{2}mv^{2}-\frac{1}{2}mv^{2}=0$，故 A 错误；从第 $3\Us$ 末到第 $5\Us$ 末，合力做的功为 $W_{2}=0-\frac{1}{2}mv^{2}=-W$，故 B 错误；从第 $5\Us$ 末到第 $7\Us$ 末，合力做的功为 $W_{2}=\frac{1}{2}mv^{2}-0=W$，故 C 正确；从第 $3\Us$ 末到第 $4\Us$ 末，合力做的功为 $W_{2}=\frac{1}{2}m\left(\frac{v}{2}\right)^{2}-\frac{1}{2}mv^{2}=-0.75W$，故 D 正确。
}

\item
%% number: 13
%% paperName: 2007年上海高考第 13 题 5 分
%% typeId: 2
%% type: 多选题
%% chapter: 电场
%% point: 带电粒子的往复运动
%% method: 无
%% score: 5
%% degree: 450
%% duplicateId: 0
%% body:
一点电荷仅受电场力作用，由 A 点无初速释放，先后经过电场中的 B 点和 C 点。点电荷在 A、B、C 三点的电势能分别用 $E_{A}$、$E_{B}$、$E_{C}$ 表示，则 $E_{A}$、$E_{B}$ 和 $E_{C}$ 间的关系可能是\xzanswer{AD}

\fourchoices[answer=AD]
{$E_{A}>E_{B}>E_{C}$}
{$E_{A}<E_{B}<E_{C}$}
{$E_{A}<E_{C}<E_{B}$}
{$E_{A}>E_{C}>E_{B}$}

%% answer: AD
%% memo:
\memoanswer{
由于在电场力作用下，电荷由 A 运动到 B、C，电荷在 B、C 的速度大于等于零，且只有动能与电势能之间的转化，由能量守恒得，B、C 两点的电势能不高于 A 点。如在匀强电场中，电荷与电场线平行运动，电势能一直减小，故 A 正确；如在运动过程中进入到一个方向相反的匀强电场中，电荷速度减小，动能减小，电势能将增加，故 D 正确。
}

\item
%% number: 14
%% paperName: 2007年上海高考第 14 题 5 分
%% typeId: 4
%% type: 实验题
%% chapter: 直线运动
%% point: 运动图象
%% method: 无
%% score: 5
%% degree: 650
%% duplicateId: 0
%% body:
在实验中得到小车做直线运动的 $s-t$ 关系如图所示。

\begin{enumerate}
	\item
	由图可以确定，小车在 AC 段和 DE 段的运动分别为（　）

	（A）AC 段是匀加速运动；DE 段是匀速运动

	（B）AC 段是加速运动；DE 段是匀加速运动

	（C）AC 段是加速运动；DE 段是匀速运动

	（D）AC 段是匀加速运动；DE 段是匀加速运动
	\item
	在与 AB、AC、AD 段对应的平均速度中，最接近小车在 A 点瞬时速度的是 \tkanswer[1.6cm]{AB} 段中的平均速度。
\end{enumerate}

\onepicture[w=6.5cm, align=right]{\includesvg{figs/14}}

%% answer:
\jdanswer{
\begin{enumerate}
	\item
	C

	\item
	AB
\end{enumerate}
}

%% memo:
\memoanswer{
（1）由 $s-t$ 关系图线知，AC 段各点的切线斜率逐渐变大，故为加速运动，不能确定是否为匀加速运动；DE 段为倾斜直线，故是匀速运动，故 C 正确。

（2）由瞬时速度的概念知 $v=\frac{\Delta s}{\Delta t}$，$\Delta t\to 0$ 时 $v$ 为物体运动到某点的瞬时速度，故最接近小车在 A 点瞬时速度的是 AB 段中的平均速度。
}

\item
%% number: 15
%% paperName: 2007年上海高考第 15 题 6 分
%% typeId: 4
%% type: 实验题
%% chapter: 电路
%% point: 电路实验
%% method: 等效替代
%% score: 6
%% degree: 650
%% duplicateId: 0
%% body:
为了测量一个阻值较大的未知电阻，某同学使用了干电池（$1.5\UV$），毫安表（$1\UmA$），电阻箱（$0\sim9999\UO$），电键，导线等器材。该同学设计的实验电路如图（a）所示，实验时，将电阻箱阻值置于最大，断开 $K_{2}$，闭合 $K_{1}$，减小电阻箱的阻值，使电流表的示数为 $I_{1}=1.00\UmA$，记录电流强度值；然后保持电阻箱阻值不变，断开 $K_{1}$，闭合 $K_{2}$，此时电流表的示数为 $I_{2}=0.80\UmA$，记录电流强度值。由此可得被测电阻的阻值为 \tkanswer[2.2cm]{375} $\UO$。经分析，该同学认为上述方案中电源电动势的值可能与标称值不一致，因此会造成误差。为避免电源对实验结果的影响，又设计了如图（b）所示的实验电路，实验过程如下：断开 $K_{1}$，闭合 $K_{2}$，此时电流表指针处于某一位置，记录相应的电流值，其大小为 $I$；断开 $K_{2}$，闭合 $K_{1}$，调节电阻箱的阻值，使电流表的示数为 \tkanswer[1.4cm]{$I$}，记录此时电阻箱的阻值，其大小为 $R_{0}$。由此可测出 $R_{x}=$ \tkanswer[1.8cm]{$R_{0}$}。

\twopicture[labelA=2007上海15a, labelB=2007上海15b, wA=5.5cm, wB=5.5cm, align=right]
{figs/15a.png}
{figs/15b.png}

%% answer:
\jdanswer{375，$I$，$R_{0}$}

%% memo:
\memoanswer{
（1）断开 $K_{2}$ 闭合 $K_{1}$，此时 $I_{1}=1.00\UmA$，电路中有电阻箱 $R_{0}$、电源内阻 $r$、电流表内阻 $R_{A}$ 三部分电阻，由闭合电路欧姆定律得 $E=I_{1}(R_{A}+R_{0}+r)$；断开 $K_{1}$ 闭合 $K_{2}$，此时 $I_{2}=0.80\UmA$，电路中有电阻箱 $R_{0}$、电源内阻 $r$、电流表内阻 $R_{A}$ 和待测电阻 $R_{x}$ 四部分电阻，由闭合电路欧姆定律得 $E=I_{2}(R_{A}+R_{0}+r+R_{x})$，将数据代入上述方程得 $1.5\UV=1.00\UmA\times(R_{A}+R_{0}+r)$、$1.5\UV=0.80\UmA\times(R_{A}+R_{0}+r+R_{x})$，由前式解得 $R_{A}+R_{0}+r=1500\UO$，代入后式可解得 $R_{x}=375\UO$。

（2）在等效替代法测电阻的实验中，两次电流表示数相等，说明两次电路中的电阻相同，即电阻箱和待测电阻的阻值等效，故调节电阻箱的阻值，使电流表的示数为 $I$，此时 $R_{x}=R_{0}$。
}

\item
%% number: 16
%% paperName: 2007年上海高考第 16 题 5 分
%% typeId: 4
%% type: 实验题
%% chapter: 电路
%% point: 输出功率极值
%% method: 极值
%% score: 5
%% degree: 450
%% duplicateId: 0
%% body:
某同学设计了如图（a）所示电路研究电源输出功率变化情况。电源 E 的电动势、内电阻恒定，$R_{1}$ 为滑动变阻器，$R_{2}$、$R_{3}$ 为定值电阻，A、V 为理想电表。

\begin{enumerate}
	\item
	若滑动片 P 由 a 滑至 b 时 A 读数一直变小，则 $R_{1}$ 和 $R_{2}$ 必须满足的关系是 \tkanswer[2.2cm]{$R_{1}\leqslant R_{2}$}。
	\item
	若 $R_{1}=6\UO$，$R_{2}=12\UO$，电源内电阻 $r=6\UO$，当滑动片 P 由 a 滑至 b 时，电源 E 的输出功率 $P$ 随外电路的总电阻 $R$ 的变化关系如图（b）所示，则 $R_{3}$ 的阻值应该选择（　）

	（A）$2\UO$　　（B）$4\UO$　　（C）$6\UO$　　（D）$8\UO$
\end{enumerate}

\twopicture[labelA=2007上海16a, labelB=2007上海16b, wA=4.0cm, wB=4.6cm, align=right]
{\includesvg{figs/16a}}
{\includesvg{figs/16b}}

%% answer:
\jdanswer{
\begin{enumerate}
	\item
	$R_{1}\leqslant R_{2}$

	\item
	B
\end{enumerate}
}

%% memo:
\memoanswer{
（1）A 示数为 $I=\frac{E}{r+R_{3}+R_{\text{并}}}$，欲使滑动片 P 由 a 滑至 b 时示数一直变小，须使 $R_{\text{并}}$ 一直变大，因为两电阻之和一定时，当二者相等时并联电阻最大，所以要求 $R_{1}\leqslant R_{2}$。

（2）从图（b）知，当 $R_{\text{外}}=6\UO$ 时，电源 E 的输出功率 $P$ 最大，即电源内阻为 $6\UO$，而 $R_{1}$ 和 $R_{2}$ 部分的电阻变化范围是 $0\sim4\UO$，因此 $R_{3}$ 的阻值必须大于 $2\UO$ 小于 $6\UO$，故 B 正确。
}

\item
%% number: 17
%% paperName: 2007年上海高考第 17 题 8 分
%% typeId: 4
%% type: 实验题
%% chapter: 曲线运动
%% point: 平抛运动
%% method: 无
%% score: 8
%% degree: 450
%% duplicateId: 0
%% body:
利用单摆验证小球平抛运动规律，设计方案如图（a）所示，在悬点 O 正下方有水平放置的炽热的电热丝 p，当悬线摆至电热丝处时能轻易被烧断；MN 为水平木板，已知悬线长为 $L$，悬点与木板间的距离 $OO'=h$（$h>L$）。

\begin{enumerate}
	\item
	电热丝 p 必须置于悬点正下方的理由是：\tkanswer[3.6cm]{保证小球沿水平方向抛出}。
	\item
	将小球向左拉起后自由释放，最后小球落在木板上的 C 点，$O'C=S$，则小球做平抛运动的初速度 $v_{0}$ 为 \tkanswer[3.0cm]{$S\sqrt{\frac{g}{2(h-L)}}$}。
	\item
	在其他条件不变的情况下，若改变释放小球时悬线与竖直方向间的夹角 $\theta$，小球落点与 $O'$ 点的水平距离 $S$ 将随之改变，经多次实验，以 $S^{2}$ 为纵坐标、$\cos\theta$ 为横坐标，得到如图（b）所示图象。则当 $\theta=30^{\circ}$ 时，$S$ 为 \tkanswer[1.2cm]{0.52} $\Um$；若悬线长 $L=1.0\Um$，悬点与木板间的距离 $OO'$ 为 \tkanswer[1.2cm]{1.5} $\Um$。
\end{enumerate}

\twopicture[labelA=2007上海17a, labelB=2007上海17b, wA=4.2cm, wB=4.8cm, align=right]
{\includesvg{figs/17a}}
{\includesvg{figs/17b}}

%% answer:
\jdanswer{
\begin{enumerate}
	\item
	保证小球沿水平方向抛出

	\item
	$S\sqrt{\frac{g}{2(h-L)}}$

	\item
	0.52，1.5
\end{enumerate}
}

%% memo:
\memoanswer{
（1）由于在烧断细线前小球做圆周运动，故速度方向沿切线方向，所以只有在悬点正下方小球的速度沿水平方向。要小球做平抛运动，则小球平抛的初速度只能沿水平方向，故只有保证小球沿水平方向抛出才能保证小球做平抛运动。

（2）由于小球做平抛运动，由平抛运动规律，在水平方向有 $S=v_{0}t$，在竖直方向有 $h-L=\frac{1}{2}gt^{2}$，解得 $v_{0}=S\sqrt{\frac{g}{2(h-L)}}$。

（3）改变释放小球时悬线与竖直方向的夹角 $\theta$ 时，小球平抛的速度 $v$，由动能定理得 $mg(L-L\cos\theta)=\frac{1}{2}mv^{2}$，在水平方向有 $S=vt$，联立可得 $S^{2}=4(h-L)L(1-\cos\theta)$。由图像知，当 $\cos\theta=0$ 时有 $2=4(h-L)L$，化简得 $S^{2}=2-2\cos\theta$；当 $\theta=30^{\circ}$ 时 $\cos\theta=\frac{\sqrt{3}}{2}$，解得 $S=\sqrt{2-\sqrt{3}}\approx0.52\Um$。当 $L=1.0\Um$ 时，$S^{2}=4(h-1)(1-\cos\theta)$，故当 $\cos\theta=0$ 时有 $2=4(h-1)$，解得 $h=1.5\Um$。
}

\item
%% number: 18
%% paperName: 2007年上海高考第 18 题 6 分
%% typeId: 4
%% type: 实验题
%% chapter: 气体性质
%% point: 气体拓展实验
%% method: 无
%% score: 6
%% degree: 450
%% duplicateId: 0
%% body:
一定量的理想气体与两种实际气体 Ⅰ、Ⅱ 在标准大气压下做等压变化时的 $V-T$ 关系如图（a）所示，图中 $\frac{V'-V_{0}}{V_{0}-V''}=\frac{1}{2}$。用三份上述理想气体作为测温物质制成三个相同的气体温度计，然后将其中二个温度计中的理想气体分别换成上述实际气体 Ⅰ、Ⅱ。在标准大气压下，当环境温度为 $T_{0}$ 时，三个温度计的读数各不相同，如图（b）所示，温度计（ii）中的测温物质应为实际气体 \tkanswer[1.2cm]{Ⅱ}（图中活塞质量忽略不计）；若此时温度计（ii）和（iii）的读数分别为 $21^{\circ}\mathrm{C}$ 和 $24^{\circ}\mathrm{C}$，则此时温度计（i）的读数为 \tkanswer[1.4cm]{23} $^{\circ}\mathrm{C}$；可见用实际气体作为测温物质时，会产生误差。为减小在 $T_{1}-T_{2}$ 范围内的测量误差，现针对 $T_{0}$ 进行修正，制成如图（c）所示的复合气体温度计，图中无摩擦导热活塞将容器分成两部分，在温度为 $T_{1}$ 时分别装入适量气体 Ⅰ 和 Ⅱ，则两种气体体积之比 $V_{\mathrm{I}}:V_{\mathrm{II}}$ 应为 \tkanswer[1.6cm]{$2:1$}。

\threepicture[labelA=2007上海18a, labelB=2007上海18b, labelC=2007上海18c, wA=4.6cm, wB=3.2cm, wC=2.2cm, align=right]
{\includesvg{figs/18a}}
{figs/18b.png}
{figs/18c.png}

%% answer:
\jdanswer{Ⅱ，23，$2:1$}

%% memo:
\memoanswer{
从 $V-T$ 图象可看出，温度为 $T_{1}$ 时，气体的体积相同；温度为 $T_{0}$ 时，Ⅱ 的体积小，所以温度计（ii）中的测温物质应为实际气体 Ⅱ。在温度为 $T_{0}$ 时，$T_{\text{ii}}=21^{\circ}\mathrm{C}$，$T_{\text{iii}}=24^{\circ}\mathrm{C}$，由于 $\frac{V'-V_{0}}{V_{0}-V''}=\frac{T_{\text{iii}}-T_{0}}{T_{0}-T_{\text{ii}}}=\frac{1}{2}$，代入数据可得 $T_{0}=23^{\circ}\mathrm{C}$。对 $T_{0}$ 校正，在温度 $T_{1}$ 时装入气体体积比为 $2:1$。
}

\item
%% number: 19
%% paperName: 2007年上海高考第 19 题 10 分
%% typeId: 6
%% type: 计算题
%% chapter: 运动和力
%% point: 牛顿定律的应用
%% method: 无
%% score: 10
%% degree: 650
%% duplicateId: 0
%% body:
\textbf{A．}宇航员在地球表面以一定初速度竖直上抛一小球，经过时间 $t$ 小球落回原处；若他在某星球表面以相同初速度竖直上抛同一小球，需经过时间 $5t$ 小球落回原处。（取地球表面重力加速度 $g=10\Umsq$，阻力不计）

\begin{enumerate}
	\item
	求该星球表面附近的重力加速度 $g'$；
	\item
	已知该星球的半径与地球半径之比为 $R_{\text{星}}:R_{\text{地}}=1:4$，求该星球的质量与地球质量之比 $M_{\text{星}}:M_{\text{地}}$。
\end{enumerate}

\textbf{B．}固定光滑细杆与地面成一定倾角，在杆上套有一个光滑小环，小环在沿杆方向的推力 $F$ 作用下向上运动，推力 $F$ 与小环速度 $v$ 随时间变化规律如图所示，取重力加速度 $g=10\Umsq$。求：

\begin{enumerate}
	\item
	小环的质量 $m$；
	\item
	细杆与地面间倾角 $\alpha$。
\end{enumerate}

\threepicture[labelA=2007上海19a, labelB=2007上海19b, labelC=2007上海19c, wA=4.0cm, wB=4.4cm, wC=4.2cm, align=right]
{figs/19a.png}
{figs/19b.png}
{figs/19c.png}

%% answer:
\jdanswer{
\textbf{A．}
\begin{enumerate}
	\item
	$g'=\frac{g}{5}=2\Umsq$

	\item
	$1:80$
\end{enumerate}

\textbf{B．}
\begin{enumerate}
	\item
	$m=1\Ukg$

	\item
	$\alpha=30^{\circ}$
\end{enumerate}
}

%% memo:
\memoanswer{
\textbf{A．}（1）设竖直上抛初速度为 $v_{0}$，则 $v_{0}=\frac{gt}{2}=g'\cdot\frac{5t}{2}$，解得 $g'=\frac{g}{5}=2\Umsq$。

（2）设小球质量为 $m$，由万有引力定律得 $mg=\frac{GMm}{R^{2}}$，解得 $M=\frac{gR^{2}}{G}$，故 $M_{\text{星}}:M_{\text{地}}=\frac{g'R_{\text{星}}^{2}}{gR_{\text{地}}^{2}}=\frac{1}{5}\times\frac{1}{16}=\frac{1}{80}=1:80$。

\textbf{B．}（1）由 $F-t$ 图知在 $0\sim2\Us$ 内，由牛顿第二定律得 $F_{1}-mg\sin\alpha=ma$ ①；由 $v-t$ 图得 $a=0.5\Umsq$。$2\Us$ 以后，$F_{2}=mg\sin\alpha$ ②。由①、②得 $F_{1}-F_{2}=ma$，解得 $m=1\Ukg$。

（2）由②式，$\alpha=\sin^{-1}\left(\frac{F_{2}}{mg}\right)=\sin^{-1}\frac{5}{10}=30^{\circ}$。
}

\item
%% number: 20
%% paperName: 2007年上海高考第 20 题 12 分
%% typeId: 6
%% type: 计算题
%% chapter: 气体性质
%% point: 气体多过程
%% method: 临界
%% score: 12
%% degree: 650
%% duplicateId: 0
%% body:
如图所示，水平放置的汽缸内壁光滑，活塞厚度不计，在 A、B 两处设有限制装置，使活塞只能在 A、B 之间运动，B 左面汽缸的容积为 $V_{0}$，A、B 之间的容积为 $0.1V_{0}$。开始时活塞在 B 处，缸内气体的压强为 $0.9p_{0}$（$p_{0}$ 为大气压强），温度为 $297\UK$，现缓慢加热缸内气体，直至 $399.3\UK$。求：

\begin{enumerate}
	\item
	活塞刚离开 B 处时的温度 $T_{B}$；
	\item
	缸内气体最后的压强 $p$；
	\item
	在右图中画出整个过程的 $p-V$ 图线。
\end{enumerate}

\twopicture[labelA=2007上海20a, labelB=2007上海20b, wA=6.5cm, wB=6.5cm, align=right]
{figs/20a.png}
{figs/20b.png}

%% answer:
\jdanswer{
\begin{enumerate}
	\item
	$T_{B}=330\UK$

	\item
	$p=1.1p_{0}$

	\item
	如图
\end{enumerate}
}

%% memo:
\memoanswer{
（1）由题意知，该变化为等容变化，由查理定律得 $\frac{0.9p_{0}}{T_{0}}=\frac{p_{0}}{T_{B}}$，解得 $T_{B}=\frac{T_{0}}{0.9}=\frac{297}{0.9}\UK=330\UK$。

（2）考虑始、末状态间的关系，设活塞最终可以移动到 A 处，由理想气体状态方程 $\frac{0.9p_{0}V_{0}}{T_{0}}=\frac{1.1pV_{0}}{T}$，解得 $p=\frac{0.9T}{1.1T_{0}}p_{0}=\frac{0.9\times399.3}{1.1\times297}p_{0}=1.1p_{0}$。由结果 $p>p_{0}$，活塞可以移动到 A 处的假设成立，故缸内气体最后的压强为 $1.1p_{0}$。

（3）由上分析知，缸内气体先做等容变化，再做等压变化，最后做等容变化，图线如图所示。

\onepicture[w=5.5cm]{figs/20_an.png}
}

\item
%% number: 21
%% paperName: 2007年上海高考第 21 题 12 分
%% typeId: 6
%% type: 计算题
%% chapter: 运动和力
%% point: 牛顿定律的应用
%% method: 无
%% score: 12
%% degree: 450
%% duplicateId: 0
%% body:
如图所示，物体从光滑斜面上的 A 点由静止开始下滑，经过 B 点后进入水平面（设经过 B 点前后速度大小不变），最后停在 C 点。每隔 $0.2\Us$ 通过速度传感器测量物体的瞬时速度，下表给出了部分测量数据（重力加速度 $g=10\Umsq$）。求：

\begin{table}[htbp]
\centering
\begin{tblr}{|c|c|c|c|c|c|c|c|}
\hline
$t$（$\Us$） & 0.0 & 0.2 & 0.4 & $\cdots$ & 1.2 & 1.4 & $\cdots$ \\
\hline
$v$（$\Ums$） & 0.0 & 1.0 & 2.0 & $\cdots$ & 1.1 & 0.7 & $\cdots$ \\
\hline
\end{tblr}
\end{table}

\begin{enumerate}
	\item
	斜面的倾角 $\alpha$；
	\item
	物体和水平面之间的动摩擦因数 $\mu$；
	\item
	$t=0.6\Us$ 时的瞬时速度 $v$。
\end{enumerate}

\onepicture[w=6.5cm, align=right]{\includesvg{figs/21}}

%% answer:
\jdanswer{
\begin{enumerate}
	\item
	$\alpha=30^{\circ}$

	\item
	$\mu=0.2$

	\item
	$v=2.3\Ums$
\end{enumerate}
}

%% memo:
\memoanswer{
（1）物体在斜面上运动时为匀加速运动，加速度的大小为 $a_{1}=g\sin\alpha$。由表格中数据知 $a_{1}=\frac{\Delta v}{\Delta t}=\frac{1-0}{0.2-0}\Umsq=5\Umsq$，代入解得 $\alpha=30^{\circ}$。

（2）物体在平面上运动时为匀减速运动，加速度的大小为 $a_{2}=-\mu g$。由表格中数据知 $a_{2}=\frac{\Delta v}{\Delta t}=\frac{0.7-1.1}{1.4-1.2}=-2\Umsq$，代入解得 $\mu=0.2$。

（3）物体运动到 B 点的速度为 $v_{B}$，运动的时间为 $t_{B}$。物体在平面上运动时，有 $\frac{1.1-v_{B}}{1.2-t_{B}}=-2$，$v_{B}-1.1=-2(t_{B}-1.2)$；物体在斜面上运动时，有 $\frac{v_{B}}{t_{B}}=5$，$v_{B}=5t_{B}$，联立解得 $t_{B}=0.5\Us$。即 $t=0.6\Us$ 时物体在平面上，设其速度为 $v$，有 $v=1.1-2(0.6-1.2)\Ums=2.3\Ums$。
}

\item
%% number: 22
%% paperName: 2007年上海高考第 22 题 13 分
%% typeId: 6
%% type: 计算题
%% chapter: 电场
%% point: 带电粒子的偏转
%% method: 临界
%% score: 13
%% degree: 450
%% duplicateId: 0
%% body:
如图所示，边长为 $L$ 的正方形区域 abcd 内存在着匀强电场。电量为 $q$、动能为 $E_{k}$ 的带电粒子从 a 点沿 ab 方向进入电场，不计重力。

\begin{enumerate}
	\item
	若粒子从 c 点离开电场，求电场强度的大小和粒子离开电场时的动能；
	\item
	若粒子离开电场时动能为 $E_{k}'$，则电场强度为多大？
\end{enumerate}

\onepicture[w=5.5cm, align=right]{figs/22.png}

%% answer:
\jdanswer{
\begin{enumerate}
	\item
	$E=\frac{4E_{k}}{qL}$，$E_{k}'=5E_{k}$

	\item
	若粒子由 bc 边离开电场区域，$E=\frac{2\sqrt{E_{k}(E_{k}'-E_{k})}}{qL}$；若粒子由 cd 边离开电场区域，$E=\frac{E_{k}'-E_{k}}{qL}$
\end{enumerate}
}

%% memo:
\memoanswer{
（1）设入射速度为 $v_{0}$，粒子由 c 点离开电场，粒子做类平抛运动，由运动学公式和牛顿第二定律得 $L=\frac{1}{2}at^{2}$ ①，$L=v_{0}t$ ②，$qE=ma$ ③。由①、②、③式得 $E=\frac{4E_{k}}{qL}$ ④。

由动能定理得 $E_{k}'=E_{k}+qEL$，以④代入得 $E_{k}'=5E_{k}$。

（2）粒子可能由 bc 也可能由 cd 边离开电场，由 c 点离开电场为临界点。当粒子由 c 点离开电场时有 $E_{k}'=5E_{k}$。

若粒子由 bc 边离开电场区域，有 $E_{k}'\leqslant5E_{k}$，由动能定理得 $E_{k}'-E_{k}=qEy$；由运动学公式得 $y=\frac{1}{2}at^{2}=\frac{1}{2}\frac{qE}{m}\left(\frac{L}{v_{0}}\right)^{2}=\frac{qEL^{2}}{4E_{k}}$，代入得 $E_{k}'-E_{k}=\frac{q^{2}E^{2}L^{2}}{4E_{k}}$，解得 $E=\frac{2\sqrt{E_{k}(E_{k}'-E_{k})}}{qL}$。

若粒子由 cd 边离开电场区域，有 $E_{k}'>5E_{k}$，由动能定理得 $E_{k}'-E_{k}=qEL$，解得 $E=\frac{E_{k}'-E_{k}}{qL}$。
}

\item
%% number: 23
%% paperName: 2007年上海高考第 23 题 13 分
%% typeId: 6
%% type: 计算题
%% chapter: 电磁感应
%% point: 电磁感应能量分析
%% method: 无
%% score: 13
%% degree: 250
%% duplicateId: 0
%% body:
如图（a）所示，光滑的平行长直金属导轨置于水平面内，间距为 $L$，导轨左端接有阻值为 $R$ 的电阻，质量为 $m$ 的导体棒垂直跨接在导轨上。导轨和导体棒的电阻均不计，且接触良好。在导轨平面上一矩形区域内存在着竖直向下的匀强磁场，磁感应强度大小为 $B$。开始时，导体棒静止于磁场区域的右端，当磁场以速度 $v_{1}$ 匀速向右移动时，导体棒随之开始运动，同时受到水平向左、大小为 $f$ 的恒定阻力，并很快达到恒定速度，此时导体棒仍处于磁场区域内。

\begin{enumerate}
	\item
	求导体棒所达到的恒定速度 $v_{2}$；
	\item
	为使导体棒能随磁场运动，阻力最大不能超过多少？
	\item
	导体棒以恒定速度运动时，单位时间内克服阻力所做的功和电路中消耗的电功率各为多大？
	\item
	若 $t=0$ 时磁场由静止开始水平向右做匀加速直线运动，经过较短时间后，导体棒也做匀加速直线运动，其 $v-t$ 关系如图（b）所示，已知在时刻 $t$ 导体棒的瞬时速度大小为 $v_{t}$，求导体棒做匀加速直线运动时的加速度大小。
\end{enumerate}

\twopicture[labelA=2007上海23a, labelB=2007上海23b, wA=5.0cm, wB=4.6cm, align=right]
{figs/23a.png}
{figs/23b.png}

%% answer:
\jdanswer{
\begin{enumerate}
	\item
	$v_{2}=v_{1}-\frac{fR}{B^{2}L^{2}}$

	\item
	$f_{m}=\frac{B^{2}L^{2}v_{1}}{R}$

	\item
	$P_{\text{导体棒}}=f\left(v_{1}-\frac{fR}{B^{2}L^{2}}\right)$，$P_{\text{电路}}=\frac{f^{2}R}{B^{2}L^{2}}$

	\item
	$a=\frac{B^{2}L^{2}v_{t}+fR}{B^{2}L^{2}t-mR}$
\end{enumerate}
}

%% memo:
\memoanswer{
（1）磁场以恒定速度 $v_{1}$ 向右移动（若导体棒不动），电阻、导体棒和导轨构成的回路内磁场区域的面积变小，则 $\frac{\Delta S_{1}}{\Delta t}=Lv_{1}$；导体棒以恒定速度 $v_{2}$ 向右移动，回路面积变大，则 $\frac{\Delta S_{2}}{\Delta t}=Lv_{2}$；两者一起运动，回路中磁场区域的面积变小，有 $\frac{\Delta S}{\Delta t}=L(v_{1}-v_{2})$。

回路中磁通量变化，产生电动势和电流 $E=\frac{\Delta\Phi}{\Delta t}=B\frac{\Delta S}{\Delta t}=BL(v_{1}-v_{2})$，$I=\frac{E}{R}=\frac{BL}{R}(v_{1}-v_{2})$。导体棒匀速运动，安培力与阻力相等，$F_{\text{安}}=BIL=\frac{B^{2}L^{2}}{R}(v_{1}-v_{2})=f$ ①，解得 $v_{2}=v_{1}-\frac{fR}{B^{2}L^{2}}$ ②。

（2）由②式，$f$ 越大，$v_{2}$ 越小，$v_{2}$ 最小为 0，有当 $v_{2}=0$ 时，解得 $f_{m}=\frac{B^{2}L^{2}v_{1}}{R}$。

（3）设导体棒以恒定速度 $v_{2}$ 运动时，有 $P_{\text{导体棒}}=fv_{2}$，以②结果代入得 $P_{\text{导体棒}}=f\left(v_{1}-\frac{fR}{B^{2}L^{2}}\right)$；电路中消耗的功率为 $P_{\text{电路}}=I^{2}R=\frac{B^{2}L^{2}}{R}(v_{1}-v_{2})^{2}$，以①、②结果代入得 $P_{\text{电路}}=f(v_{1}-v_{2})=\frac{f^{2}R}{B^{2}L^{2}}$。

（4）设磁场的加速度为 $a$，金属棒的加速度为 $a'$，当金属棒以一定速度 $v$ 运动时，受安培力和阻力作用，由牛顿第二定律得 $\frac{B^{2}L^{2}}{R}(at-v)-f=ma'$。由图可知，在 $t$ 时刻导体棒的瞬时速度大小为 $v_{t}$，此时棒做匀加速运动，磁场与棒之间速度差为恒量，因此必有 $a=a'$，即 $\frac{B^{2}L^{2}}{R}(at-v_{t})-f=ma$，联立解得 $a'=a=\frac{B^{2}L^{2}v_{t}+fR}{B^{2}L^{2}t-mR}$。
}

\end{enumerate}

\end{document}
